\section{Compressed inventories of normal bundle midsections}
\label{sec:prismatic-inventory}

The classifier of Section~\ref{sec:cut-products} counts wholly prismatic
pieces but does not identify their normal bases. We now recover a binary
inventory of those midsections with multiplicities. Marked core components
remain counted, without a claimed description of their topology or geometry.
This is not the inventory of the full parallelity bundle attached inside
a core component.

\subsection{Additive chamber weights identify actual midsections}

Use the source-owned chamber system from Section~\ref{sec:cut-complement}.
Give each exceptional chamber weight one in a leading marker coordinate.
Give each unmarked triangle-prism chamber one in its original triangle
coordinate, and each unmarked quad-prism chamber one in its original quad
coordinate. The total dimension is $1+7T$. The quad interiors are one
interval per occupied type, and triangle-arm interiors are one interval
per arm. Along with the at most $6T$ singleton marks, there are at most
$11T$ additive indicator intervals. None of these weights expands a
represented chamber family.

An orbit's leading sum is positive exactly when its component touches a
mark. Such rows contribute only to the core count. A zero-leading row
records one normal disc per local prism in that component; the remaining
$7T$ entries are its connected normal midsection vector. Equal vectors
are grouped with a binary orbit multiplicity. The weighted source and
its independent orbit replay determine membership; merely summing proposed
vectors to a plausible total would not establish an inventory.

\begin{proposition}[Exact binary bundle-base inventory]
The weighted chamber histogram determines the complete normal midsection
inventory of the wholly prismatic cut components, and the conservative
core count. With maximum source-coordinate bit length $B$, the returned
inventory has $O(T^2(B+\log(T+1)))$ bits. Construction and discovery retain
the polynomial encoded-input bound of the weighted orbit operation;
independent verification also charges its supplied trace length.
\end{proposition}
\begin{proof}
Each chamber receives either a marker or exactly one normal-disc indicator.
Therefore a zero-leading orbit consists entirely of prisms. By the preceding
midsection proposition, its indicators are the normal triangles and
quadrilaterals of one connected midsection. Summing them gives its full
normal vector. All orbits with that vector contribute to its multiplicity.
Positive-leading orbits are exactly the conservative core components.

There are at most $v+2q\le6T$ distinct midsection vectors by
Section~\ref{sec:cut-products}. Each entry is bounded by the corresponding
cutting coordinate: the total number of prism chambers of a given local
type is at most its source disc count. Multiplicities are bounded by the
$T+D$ chamber universe, and Euler values have polynomially bounded digits.
This proves the stated output size. The weight dimension and indicator
count are polynomial in $T$, their coefficients are zero or one, and their
endpoints use $O(B+\log(T+1))$ bits. The existing weighted interval operation
has its polynomial encoded-input cost\cite[Theorem 16]{aht-orbits}, with
all dimension, run transport, arithmetic and trace work charged.
\end{proof}

The full normal coordinates are retained rather than divided by their gcd.
A connected annular midsection can be the normal double of a one-sided
M\"obius band. Dividing its vector by two changes the base and its horizontal
boundary, even when the two associated total interval bundles happen to
share an unpatterned homeomorphism type. For $m$ sheets of the retained
M\"obius vector, an even $m$ produces one M\"obius base and $m/2-1$ annular
bases, while odd $m$ produces $(m-1)/2$ annular bases. These are grouped
normal vectors, not expanded surface lists.

\subsection{Independent weights and trace reuse}

The native API is \code{normal\_prismatic\_inventory}. It validates the
actual source and coordinates, constructs the chamber intervals and
weights, and applies the existing weighted orbit kernel. Its result
contains core and prismatic counts, full normal base coordinates,
multiplicities and independently computable Euler values. It does not
assert a knot verdict.

The checker reconstructs the chamber graph through the independent face
segment model. Its own block-by-block weight construction identifies
endpoint marks and prism interiors; it imports neither weight nor inventory
producer. After exact weighted replay, it reconstructs every proposed base,
checks normal matching, source coordinate bounds and the type bound, and
compares the entire ordered inventory. Altered weights, base coordinates,
multiplicities, Euler values or source identity do not authenticate a result.
The schema is \code{normal-prismatic-inventory-v1}; signed hexadecimal
integer transport is retained.

A caller may supply a previously obtained ordinary cut-orbit certificate.
The producer then independently checks its exact fresh geometric pairings
and uses it for weighted replay, without another orbit search. A discovery
cycle cap cannot be combined with this replay path. Reuse is not a source
cache or an assertion that two inputs have the same geometry: only a trace
that matches the current owned interval system is accepted, and the inventory
proof binds the current source and vector separately.

Cycle limits constrain the underlying orbit search, not all weighted
arithmetic. Callbacks apply throughout validation, construction, transport
and verification. An incomplete query has no inventory or certificate;
a replay operation limit means unverified, never a geometric conclusion.

\subsection{The remaining core problem}

The inventory describes whole cut components made entirely of prisms. It
omits the parallelity regions that attach to exceptional geometry inside
a marked core component. In particular, connected primitive cuts may have
only a core component even while containing extensive interior parallelity.
The printed parallelity-bundle algorithms address that richer construction;
we have not implemented their full core and attachment output here.

Core geometry, boundary patterns, ordered vertical attachments, arbitrary
later hierarchy-piece validation and bounded hierarchy depth remain open.
The new standalone operator is not called by recognition and changes neither
its defaults nor its general complexity bound. It provides a verified
compressed inventory prerequisite without substituting it for full cutting.

\subsection{Source inventories and complete-query measurements}
All 1,564 maintained tests pass in 301.283 seconds; the 24 focused tests
pass in 0.734 seconds. The nine new regressions check parallel meridians,
vertex links, layered sources, odd/even one-sided families, connected normal
doubles, source/weight/base/count mutations, shared trace ownership, caps,
operation limits and cancellation. Binary JSON and a 16,385-bit multiplicity
pass independent replay without expanding components.

The 68.207-second source audit computes inventories on all 5,241 saved
cut-surface instances on 48 triangulations. Every core count, full normal
base vector, multiplicity and Euler value agrees with explicit finite chamber
recovery using the independent geometric pairing construction. The inventories
account for 7,838 midsection instances. This is agreement of complete
coordinate inventories, not just agreement of total component counts.
Representative bases on all 48 sources receive fresh Regina connectivity
and Euler checks. The earlier geometric pilot supplies the full finite
midsection/cut-signature context; no unique 3-manifold classification is
inferred from those signatures alone.

All native proofs replay with weight, inventory, chamber and orbit producers
disabled. Six large cut-orbit traces are then reused with orbit search
disabled; their inventories agree exactly with fresh queries. Nine large
binary cases include one-, four- and sixteen-tetrahedron meridian families
at $2^{128}$ or $2^{500}$ scale, the $2^{16384}$ one-tetrahedron case, and
large odd/even M\"obius stacks. The meridian families have one base vector
with multiplicity $m-1$. Even M\"obius stacks retain both the one-sided
vector and its connected annular double where present; odd stacks retain
the annular vector. The audit preserves 57 representative complete source
proofs.

For the large meridian families, maximum observed weight-run counts are
six, 24 and 96 for one, four and sixteen source tetrahedra respectively;
they do not grow between the tested $2^{128}$ and $2^{500}$ multiplicities.
These are finite diagnostic observations, not the proof of the polynomial
bound. That bound comes from the source weight construction, normal-type
accounting and the existing weighted interval operation.

Five paired timing rounds measure 80 completed inventory calls and sixteen
warmups, with identical-baseline A/A arms. Both methods return the same full
midsection vectors and multiplicities. The explicit reference expands the
finite chamber graph, computes its connected components and accumulates
normal-disc counts for each unmarked component. The native arm constructs
compressed indicators, runs weighted transport, independently replays the
complete source proof and serializes its result. All source construction,
validation, reconstruction and serialization are charged. This is an explicit
graph baseline, not an optimized published parallelity-bundle constructor.
\begin{center}
\begin{tabular}{rrrrrr}
\toprule
copies & explicit ms & native ms & ratio & old A/A & new A/A\\
\midrule
16 & 0.168 & 0.880 & 0.187 & 0.982 & 1.000\\
256 & 0.575 & 0.897 & 0.641 & 0.995 & 1.000\\
4096 & 8.060 & 0.928 & 8.647 & 1.011 & 0.989\\
65536 & 154.142 & 1.033 & 148.591 & 0.991 & 1.073\\
\bottomrule
\end{tabular}
\end{center}


At 65,536 copies the paired explicit/native ratio is 148.591, with median
154.142 versus 1.033 milliseconds. At 4,096 copies the ratio is 8.647.
At sixteen and 256 copies the ratios are 0.187 and 0.641: native proof and
weighted reconstruction lose on small inputs. Old A/A ratios range from
0.982 to 1.011 and new ones from 0.989 to 1.073. These remain host observations,
not universal speedup factors or full-recognition timings.

Runtime release \code{a10e8eb83}, all tests, exact source inventories,
search-free trace reuse checks and complete-query timings are retained.

The reproduction driver is \code{normal\_orbit\_research.prismatic\_inventory}.
Records are under \code{data/prismatic-inventory-*}. Publication review checks
664 current source/evidence hashes, the exact pilot input, every timing outcome
and the passing test records. It independently replays all 57 saved source
proofs with producers disabled. The complete recognizer is unchanged. The
next geometric task is parallelity inside the marked cores and its attachment
representation; that richer problem is not replaced by this whole-component
inventory.

